% --- 题目 ---
\problem[P215-13 (23-1)]{设二维随机变量 $\displaystyle (X, Y)$ 的概率密度为
\[ f(x, y) = \begin{cases} \displaystyle \frac{2}{\pi}(x^2 + y^2), & \displaystyle x^2 + y^2 \le 1, \\ 0, & \text{其他}. \end{cases} \]
\begin{enumerate}
\item[(1)] 求 $\displaystyle X$ 与 $\displaystyle Y$ 的协方差;
\item[(2)] $\displaystyle X$ 与 $\displaystyle Y$ 是否相互独立?
\item[(3)] 求 $\displaystyle Z = X^2 + Y^2$ 的概率密度.
\end{enumerate}}
\begin{note}
  主要错的是 (2)
\end{note}
\ansat{365；【十年真题】 - 考点二：随机变量的协方差与相关系数 - 13}

% --- 题目 ---
\problem[P229-2 (22-1,3)]{设 $\displaystyle X_1, X_2, \dots, X_n$ 为来自均值为 $\displaystyle \theta$ 的指数分布总体的简单随机样本, $\displaystyle Y_1, Y_2, \dots, Y_m$ 为来自均值为 $\displaystyle 2\theta$ 的指数分布总体的简单随机样本, 且两样本相互独立, 其中 $\displaystyle \theta (\theta > 0)$ 为未知参数. 利用样本 $\displaystyle X_1, X_2, \dots, X_n, Y_1, Y_2, \dots, Y_m$, 求 $\displaystyle \theta$ 的最大似然估计量 $\displaystyle \hat{\theta}$, 并求 $\displaystyle D(\hat{\theta})$.
}
\ansat{371；【十年真题】 - 考点一：矩估计与最大似然估计 - 2}

% --- 题目 ---
\problem[P219-12 (15-1,3)]{设随机变量 $\displaystyle X$ 的概率密度为
\[ f(x) = \begin{cases} 2^{-x} \ln 2, & x > 0, \\ 0, & x \le 0. \end{cases} \]
对 $\displaystyle X$ 进行独立重复的观测, 直到第 2 个大于 3 的观测值出现时停止, 记 $\displaystyle Y$ 为观测次数.
\begin{enumerate}
\item[(1)] 求 $\displaystyle Y$ 的概率分布;
\item[(2)] 求 $\displaystyle EY$.
\end{enumerate}}
\begin{note}
  主要错的是 (2)
\end{note}
\ansat{366；【真题精选】 - 考点一：随机变量的数学期望与方差 - 12}

% --- 题目 ---
\problem[P198-11 (97-1)]{袋中有 $\displaystyle 50$ 个乒乓球, 其中 $\displaystyle 20$ 个是黄球, $\displaystyle 30$ 个是白球, 今有两人依次随机地从袋中各取一球, 取后不放回, 则第二个人取得黄球的概率是 \underline{\hspace{4em}}.}
\ansat{354；【真题精选】 - 考点一：概率的五大公式 - 11}

% --- 题目 ---
\problem[P215-10 (23-3)]{设随机变量 $\displaystyle X$ 与 $\displaystyle Y$ 相互独立, 且 $\displaystyle X \sim B(1, p)$, $\displaystyle Y \sim B(2, p), p \in (0, 1)$, 则 $\displaystyle X+Y$ 与 $\displaystyle X-Y$ 的相关系数为 \underline{\hspace{4em}}.}
\ansat{364；【十年真题】 - 考点二：随机变量的协方差与相关系数 - 10}

% --- 题目 ---
\problem[P215-8 (20-3)]{设随机变量 $\displaystyle (X, Y)$ 服从二维正态分布 $\displaystyle N\left(0, 0; 1, 4; -\frac{1}{2}\right)$, 则下列随机变量中服从标准正态分布且与 $\displaystyle X$ 独立的是~(~\quad~)}
\begin{tasks}(4)
  \task $\displaystyle \frac{\sqrt{5}}{5}(X+Y).$
  \task $\displaystyle \frac{\sqrt{5}}{5}(X-Y).$
  \task $\displaystyle \frac{\sqrt{3}}{3}(X+Y).$
  \task $\displaystyle \frac{\sqrt{3}}{3}(X-Y).$
\end{tasks}
\ansat{364；【十年真题】 - 考点二：随机变量的协方差与相关系数 - 8}

% --- 题目 ---
\problem[P222-1 (22-1)]{设随机变量 $\displaystyle X_1, X_2, \dots, X_n$ 独立同分布. 且 $\displaystyle X_1$ 的 $\displaystyle 4$ 阶矩存在. 记 $\displaystyle \mu_k = E(X_1^k) \ (k=1,2,3,4)$, 则根据切比雪夫不等式, 对任意 $\displaystyle \varepsilon > 0$, 有
\[ P\left\{\left|\frac{1}{n}\sum_{i=1}^n X_i^2 - \mu_2\right| \ge \varepsilon\right\} \le (\quad) \]
}
\begin{tasks}(4)
  \task $\displaystyle \frac{\mu_4 - \mu_2^2}{n\varepsilon^2}$.
  \task $\displaystyle \frac{\mu_4 - \mu_2^2}{\sqrt{n}\varepsilon^2}$.
  \task $\displaystyle \frac{\mu_2 - \mu_1^2}{n\varepsilon^2}$.
  \task $\displaystyle \frac{\mu_2 - \mu_1^2}{\sqrt{n}\varepsilon^2}$.
\end{tasks}
\ansat{368；【十年真题】 - 考点：大数定律、中心极限定理与切比雪夫不等式 - 1}

% --- 题目 ---
\problem[P196-7 (25-1)]{设 $\displaystyle A, B$ 为两个随机事件, 且 $\displaystyle A$ 与 $\displaystyle B$ 相互独立. 已知 $\displaystyle P(A)=2P(B)$, $\displaystyle P(A \cup B)=\frac{5}{8}$, 则在事件 $\displaystyle A, B$ 至少有一个发生的条件下, $\displaystyle A, B$ 中恰有一个发生的概率为 \underline{\hspace{4em}}.}
\ansat{353；【十年真题】 - 考点一：概率的五大公式 - 7}

% --- 题目 ---
\problem[P202-例3 (90-1)]{已知随机变量 $\displaystyle X$ 的概率密度函数 $\displaystyle f(x) = \frac{1}{2}\mathrm{e}^{-|x|}, -\infty < x < +\infty$, 则 $\displaystyle X$ 的分布函数 $\displaystyle F(x) = \underline{\hspace{4em}}. $}
\ansat{202；【方法探究】 - 考点一：随机变量的分布 - 例3}
\vspace{2em}

% --- 题目 ---
\problem[P210-4 (10-1,3)]{设二维随机变量 $\displaystyle (X, Y)$ 的概率密度为 $\displaystyle f(x, y) = A\mathrm{e}^{-2x^2 + 2xy - y^2}, -\infty < x < +\infty, -\infty < y < +\infty$, 求常数 $\displaystyle A$ 及条件概率密度 $\displaystyle f_{Y|X}(y|x)$.}
\ansat{360；【真题精选】 - 考点一：二维随机变量的分布 - 4}

% --- 题目 ---
\problem[P206-1 (24-1,3)]{设随机变量 $\displaystyle X, Y$ 相互独立, 且均服从参数为 $\displaystyle \lambda$ 的指数分布. 令 $\displaystyle Z = |X - Y|$, 则下列随机变量中与 $\displaystyle Z$ 同分布的是~(~\quad~)}
\begin{tasks}(4)
  \task $\displaystyle X + Y.$
  \task $\displaystyle \frac{X+Y}{2}.$
  \task $\displaystyle 2X.$
  \task $\displaystyle X.$
\end{tasks}
\ansat{358；【十年真题】 - 考点二：两个随机变量的函数的分布 - 1}

% --- 题目 ---
\problem[P215-12 (20-1)]{设 $\displaystyle X$ 服从区间 $\displaystyle \left(-\frac{\pi}{2}, \frac{\pi}{2}\right)$ 上的均匀分布, $\displaystyle Y = \sin X$, 则 $\displaystyle \mathrm{Cov}(X, Y) = \underline{\hspace{4em}}. $}
\ansat{365；【十年真题】 - 考点二：随机变量的协方差与相关系数 - 12}

% --- 题目 ---
\problem[P224-2 (25-3-仅数学三)]{设总体 $\displaystyle X$ 的分布函数为 $\displaystyle F(x)$, $\displaystyle X_1, X_2, \dots, X_n$ 为来自总体 $\displaystyle X$ 的简单随机样本, 样本的经验分布函数为 $\displaystyle F_n(x)$. 对于给定的 $\displaystyle x$ ($\displaystyle 0 < F(x) < 1$), $\displaystyle D[F_n(x)] = ( \quad )$
\begin{tasks}(2)
  \task $\displaystyle F(x)[1-F(x)]$.
  \task $\displaystyle [F(x)]^2$.
  \task $\displaystyle \frac{1}{n} F(x)[1-F(x)]$.
  \task $\displaystyle \frac{1}{n} [F(x)]^2$.
\end{tasks}}
\ansat{369；【十年真题】 - 考点一：抽样分布 - 2}

% --- 题目 ---
\problem[P199-14 (89-1)]{甲、乙两人独立地对同一目标射击一次, 其命中率分别为 $\displaystyle 0.6$ 和 $\displaystyle 0.5$. 现已知目标被命中, 则它是甲射中的概率为 \underline{\hspace{4em}}.}
\ansat{355；【真题精选】 - 考点一：概率的五大公式 - 14}
\vspace{6em}

% --- 题目 ---
\problem[P211-5 (09-1,3)]{袋中有 $\displaystyle 1$ 个红球, $\displaystyle 2$ 个黑球与 $\displaystyle 3$ 个白球, 现有放回地从袋中取两次, 每次取一个球. 以 $\displaystyle X, Y, Z$ 分别表示两次取球所取得的红球、黑球与白球的个数.
\begin{enumerate}
\item[(1)] 求 $\displaystyle P\{X=1|Z=0\}$;
\item[(2)] 求二维随机变量 $\displaystyle (X,Y)$ 的概率分布.
\end{enumerate}}
\ansat{360；【真题精选】 - 考点一：二维随机变量的分布 - 5}

% --- 题目 ---
\problem[P227-5 (08-1,3)]{设 $\displaystyle X_1, X_2, \dots, X_n$ 是总体 $\displaystyle N(\mu, \sigma^2)$ 的简单随机样本. 记 $\displaystyle \bar{X} = \frac{1}{n} \sum_{i=1}^n X_i, S^2 = \frac{1}{n-1} \sum_{i=1}^n (X_i - \bar{X})^2, \ T = \bar{X}^2 - \frac{1}{n} S^2$.
\begin{enumerate}
\item[(1)] 证明 $\displaystyle ET = \mu^2$;
\item[(2)] 当 $\displaystyle \mu=0, \sigma=1$ 时, 求 $\displaystyle DT$.
\end{enumerate}}
\begin{note}
  主要错的是 (2)
\end{note}
\ansat{371；【真题精选】 - 考点二：统计量的数字特征 - 5}

% --- 题目 ---
\problem[P215-5 (22-1)]{设随机变量 $\displaystyle X \sim N(0, 1)$, 在 $\displaystyle X=x$ 的条件下随机变量 $\displaystyle Y \sim N(x, 1)$, 则 $\displaystyle X$ 与 $\displaystyle Y$ 的相关系数为~(~\quad~)}
\begin{tasks}(4)
  \task $\displaystyle \frac{1}{4}.$
  \task $\displaystyle \frac{1}{2}.$
  \task $\displaystyle \frac{\sqrt{3}}{3}.$
  \task $\displaystyle \frac{\sqrt{2}}{2}.$
\end{tasks}
\ansat{364；【十年真题】 - 考点二：随机变量的协方差与相关系数 - 5}

% --- 题目 ---
\problem[P229-3 (20-1,3)]{设某种元件的使用寿命 $\displaystyle T$ 的分布函数为
\[ F(t) = \begin{cases} 1 - \mathrm{e}^{-(\frac{t}{\theta})^m}, & t \ge 0, \\ 0, & \text{其他}, \end{cases} \]
其中 $\displaystyle \theta, m$ 为参数且大于零.
\begin{enumerate}
\item[(1)] 求概率 $\displaystyle P\{T>t\}$ 与 $\displaystyle P\{T>t+s | T>s\}$, 其中 $\displaystyle s>0, t>0$;
\item[(2)] 任取 $\displaystyle n$ 个这种元件做寿命试验, 测得它们的寿命分别为 $\displaystyle t_1, t_2, \dots, t_n$. 若 $\displaystyle m$ 已知, 求 $\displaystyle \theta$ 的最大似然估计值 $\displaystyle \hat{\theta}$.
\end{enumerate}}
\begin{note}
  主要错的是 (2)
\end{note}
\ansat{371；【十年真题】 - 考点一：矩估计与最大似然估计 - 3}

% --- 题目 ---
\problem[P214-8 (25-1,3)]{投保人的损失事件发生时, 保险公司的赔付额 $\displaystyle Y$ 与投保人的损失额 $\displaystyle X$ 的关系为
\[ Y = \begin{cases} 0, & X \le 100, \\ X - 100, & X > 100. \end{cases} \]
设损失事件发生时, 投保人的损失额 $\displaystyle X$ 的概率密度为
\[ f(x) = \begin{cases} \displaystyle \frac{2 \times 100^2}{(100+x)^3}, & x > 0, \\ 0, & x \le 0. \end{cases} \]
\begin{enumerate}
\item[(1)] 求 $\displaystyle P\{Y > 0\}$ 及 $\displaystyle EY$;
\item[(2)] 这种损失事件在一年内发生的次数记为 $\displaystyle N$, 保险公司在一年内就这种损失事件产生的理赔次数记为 $\displaystyle M$. 假设 $\displaystyle N$ 服从参数为 8 的泊松分布, 在 $\displaystyle N=n \ (n \ge 1)$ 的条件下, $\displaystyle M$ 服从二项分布 $\displaystyle B(n, p)$, 其中 $\displaystyle p=P\{Y > 0\}$. 求 $\displaystyle M$ 的概率分布.
\end{enumerate}}
\begin{note}
  主要错的是 (2)
\end{note}
\ansat{363；【十年真题】 - 考点一：随机变量的数学期望与方差 - 8}

% --- 题目 ---
\problem[P204-3 (13-1)]{设随机变量 $\displaystyle X$ 的概率密度为 $\displaystyle f(x) = \begin{cases} \frac{1}{9}x^2, & 0 < x < 3, \\ 0, & \text{其他}. \end{cases}$, 
\newline
令随机变量 $\displaystyle Y = \begin{cases} 2, & X \le 1, \\ X, & 1 < X < 2, \\ 1, & X \ge 2. \end{cases}.$
\begin{enumerate}
\item[(1)] 求 \(Y\) 的分布函数;
\item[(2)] 求概率 \(P\left\{ X \leq Y \right\}\).
\end{enumerate}}
\ansat{357；【真题精选】 - 考点二：随机变量的函数的分布 - 3}

% --- 题目 ---
\problem[P196-11 (18-3)]{设随机事件 $\displaystyle A, B, C$ 相互独立, 且}
\[ P(A) = P(B) = P(C) = \frac{1}{2}, \]
{则 $\displaystyle P(AC|A \cup B) = \underline{\hspace{4em}}.$}
\ansat{354；【十年真题】 - 考点一：概率的五大公式 - 11}

% --- 题目 ---
\problem[P200-1 (23-3-局部)]{设随机变量 $\displaystyle X$ 的概率密度为 $\displaystyle f(x) = \frac{\mathrm{e}^x}{(1+\mathrm{e}^x)^2}, -\infty < x < +\infty$, 令 $\displaystyle Y = \mathrm{e}^X$.}
\begin{enumerate}
\item[(1)] 求 $\displaystyle X$ 的分布函数;
\item[(2)] 求 $\displaystyle Y$ 的概率密度.
\end{enumerate}
\begin{note}
  主要错的是 (2)
\end{note}
\ansat{355；【十年真题】 - 考点二：随机变量的函数的分布 - 1}

% --- 题目 ---
\problem[P224-3 (24-3-仅数学三)]{设总体 $\displaystyle X$ 服从 $\displaystyle [0, \theta]$ 上的均匀分布, 其中 $\displaystyle \theta \in (0, +\infty)$ 为未知参数. $\displaystyle X_1, X_2, \dots, X_n$ 是来自总体 $\displaystyle X$ 的简单随机样本, 记 $\displaystyle X_{(n)} = \max\{X_1, X_2, \dots, X_n\}, T_c = cX_{(n)}$.
\begin{enumerate}
\item[(1)] 求 $\displaystyle c$, 使得 $\displaystyle E(T_c) = \theta$;
\item[(2)] 记 $\displaystyle h(c) = E[(T_c - \theta)^2]$, 求 $\displaystyle c$ 使得 $\displaystyle h(c)$ 最小.
\end{enumerate}}
\ansat{370；【十年真题】 - 考点二：统计量的数字特征 - 3}
